\documentclass[11pt]{article}
\usepackage[margin=1.05in]{geometry}
\usepackage{booktabs}
\usepackage{amsmath}
\usepackage[T1]{fontenc}
\usepackage{lmodern}
\usepackage{microtype}
\usepackage[hidelinks]{hyperref}
\newcommand{\code}[1]{\texttt{#1}}
\title{\textbf{fmix Phase A: Structure, Rationale,\\ and the Generation Benchmarks}}
\date{2026-07-21}
\author{}
\begin{document}
\maketitle
\vspace{-3em}
\begin{center}\small documents the phase-A machinery as of commits
\code{e38d4ac0}, \code{8cb83033}, \code{8748caea}, \code{b2f4d702}
on \code{ssg-gen-mix-clean}\end{center}

\section{Context: the two-phase mixing pipeline}
Production mixing of a gadgetized circuit runs in two phases:
\begin{itemize}
\item \textbf{Phase A} (this document): churn at (roughly) constant, small
size. The input is a gadgetized circuit; the goal is to erase the
gadgetizer's deterministic authorship --- re-encode every local
neighborhood, rotate the sharing frame, homogenize the syntactic texture ---
while the circuit is small and every operation is cheap.
\item \textbf{Phase B}: anneal the DB channel away and grow by split
transport to the final size. Fossil erosion scales with the pure-split
growth ratio (R-law: ratio $\geq 3$--$4$ for $\sim$1\%-class residuals), so
phase B wants to \emph{start small} --- which is phase A's other job:
deliver its dose without inflating the circuit first.
\end{itemize}
The tension that shapes everything below: the moves that re-encode most
effectively also grow the circuit, and growth belongs to phase B. Phase A is
therefore built around a \textbf{dose accounting} (how much re-encoding has
each gate actually received?) and a \textbf{payment discipline} (spend
growth only where it is provably the only way to buy the dose).

\section{What phase A must accomplish}
Three re-randomizations, each with its own meter:
\begin{enumerate}
\item \textbf{Local nonlinear re-encoding} --- DB replacements substitute a
window (2--5 gates) with a random equivalent implementation drawn from the
frozen store. A single replacement is invisible outside its $\sim$6-wire
span; security comes from \emph{composition of overlapping rewrites}. This
is the only lever against degree-2 leakage of the original computation
(affine twists provably cannot touch it), which makes it the
security-critical axis. Its meter is the \textbf{generation census} of
\S\ref{sec:gen}.
\item \textbf{Frame erasure} --- conjugation twists rotate the wire frame
position-by-position. Meter: per-position twist coverage
$\mathrm{cov} = \text{twisted span}/\text{size}$, with the measured
saturation law $256\,(1-e^{-c/256})$ giving a target of
$\mathrm{cov}\approx 600\times$. Coverage per twist scales inversely with
circuit size --- the arithmetic reason the twist budget belongs in phase A.
\item \textbf{Texture homogenization} --- the walk's split/merge churn plus
DB re-spelling converts gadget-template material into store-native
equilibrium material and erodes origin fossils, serving the static
(non-executing) distinguisher front.
\end{enumerate}

\section{The round scheduler}
Each move is one round. Top-level coins, in order:

\begin{center}\small
\begin{tabular}{lll}
\toprule
coin & round type & size effect \\
\midrule
\code{p\_twist} & one conjugation twist (type from \code{w\_twist\_*} ratios) & $+2$/$+6$ brackets \\
\code{p\_db\_ingest} & \textbf{cheap ingest}: Compressing-mode DB attempt, cheap-tier seed & never grows \\
\code{p\_db\_hard} & \textbf{paid}: MinGrow-mode DB attempt, hard-tier seed & minimum spelling; ledgered \\
\code{p\_db\_eff} & generic size-agnostic DB round (steered/annealed) & any \\
otherwise & the walk: contract (undo/merge/\code{w\_db}) or expand (cross/split/insert) & $\pm$ small \\
\bottomrule
\end{tabular}
\end{center}

The thermostat holds size near \code{-{}-target-size}; the steer multiplies
the generic DB rate by $\sigma(-\text{excess}/\text{temp})$. An important
measured consequence: a run sitting at or above target has its generic
agnostic channel throttled to $p\approx 0.001$ --- in that regime
essentially all re-encoding flows through the non-growing and paid
channels, which is exactly the discipline phase A wants.

\section{Generation accounting}\label{sec:gen}
Every gate carries a \textbf{rewrite generation}:
\begin{itemize}
\item Input gates start at generation 0.
\item A \textbf{DB replacement} stamps its products with the \emph{upper
median} of the outgoing window's generations, plus 1 (median rounded up on
even window sizes; \code{-{}-gen-median-low} selects the lower median).
\item A \textbf{split} (presplit, cross piece, fresh-split, unsubsume,
twist case-split) stamps children with parent $+\,1$ under the default
\emph{ratchet} rule; \code{-{}-gen-split-inherit} makes children inherit the
parent generation unchanged, so that only DB replacements raise
generations.
\item \textbf{Twist bracket packets} and \textbf{fresh insert pairs} are
born-random material carrying no input structure: they get MAXGEN, higher
than every real generation.
\item Merges take the minimum of their parents; undos restore the recorded
pre-split generations.
\end{itemize}
The \textbf{circuit generation} is the largest $G$ such that at least 95\%
of \emph{all} gates have generation $\geq G$ (the 5th-percentile gate
generation, reported \code{G=}). The \textbf{dose stop} ends the run at the
first report point where the below-target fraction over all gates is
$\leq$ \code{-{}-gen-stop-frac} (0.05 = ``the circuit has reached generation
$G$'') and the twist coverage target (\code{-{}-twist-cov-stop}, if set) is
met. The move budget becomes a ceiling: phase A runs exactly as long as the
dose requires, which is the minimal-growth schedule.

\section{Targeting: the laggard census and tiers}
Uniform window seeding is a coupon-collector process --- the last
un-encoded gates absorb most of the moves. Instead, fmix maintains a census
of \textbf{laggards} (gates below \code{-{}-gen-target}), rebuilt every
\code{-{}-gen-rescan} moves and pruned lazily at draw time, and seeds DB
windows from it with probability \code{-{}-gen-bias}.

Each gate also carries a \textbf{miss counter} (reset on every splice
product): laggard-seeded attempts that fail to consume their seed bump it.
The counter partitions laggards into:
\begin{itemize}
\item \textbf{cheap tier} (\code{miss} $<$ \code{-{}-gen-miss-budget}): still
worth trying to ingest with non-growing replacements;
\item \textbf{hard tier}: the cheap channel has \emph{proven} the gate has
no non-growing spelling in its current context --- only the paid channel
touches it;
\item \textbf{retired} (\code{miss} $\geq$ \code{-{}-gen-giveup}, optional):
written off, reported \code{u=}, dropped from targeting (but still counted
by the all-gates dose criterion).
\end{itemize}

\section{Ingest-then-pay}
The policy (``first ingest as many gates as possible with tight control
over growth, then pay in growth only for those gates that cannot be
ingested otherwise'') rests on a measured dichotomy:
\begin{itemize}
\item \textbf{Compress-only regimes leave a large stuck population.} With
re-encoding restricted to non-growing replacements, $\sim$25--35\% of gates
were never consumed by any splice --- the type join was unambiguous: g57
gates 0.2\% stuck, CNOTs 89\%, non-g57 conjunctions 91\%. A CNOT alone
costs six g57s: windows containing one essentially never have a non-growing
spelling.
\item \textbf{Unrestricted any-length replacement reaches everything but
overpays.} With the agnostic channel live, the stuck set collapsed to
3.5\% --- at the cost of $4.5\times$ growth.
\end{itemize}
Ingest-then-pay separates the two channels and gates the expensive one
per-gate: \code{-{}-p-db-ingest} rounds (run hot, $\sim$0.5) use Compressing
mode --- non-growing replacements only, zero size risk --- seeded on the
cheap tier; \code{-{}-p-db-hard} rounds (low rate, $\sim$0.05) use
\textbf{MinGrow} mode --- uniform among the \emph{shortest} equivalents,
growing allowed --- seeded on the hard tier only, never falling back to
easy material, with every added gate ledgered (\code{paid=}).

The validating pilot (raw CNOT-gadgetized cg1, 48{,}165 gates, generation
target 2): dose stop fired at 900k moves with $2.12\times$ growth versus
$4.5\times$ for unrestricted agnostic mixing at the same dose;
$\code{paid}=52{,}089$ accounted for essentially all growth; the paid
channel hit on 94\% of its attempts; exactly one gate in the run was
retired as unreachable. The store can spell virtually everything --- the
earlier ``unreachable population'' was entirely the non-growing constraint,
and the miss-budget machinery locates the gates that need payment and pays
each one at its minimum spelling.

\section{The generation benchmarks}
\textbf{Setup.} Input: a gadgetized sliced sandwich on $n=128$ (two-sided
keyed slicing on 256 wires, CNOT-gadgetized to 512 wires): 99{,}016 gates.
Recipe: \code{target-size 100k}, \code{p-db-ingest 0.5},
\code{p-db-hard 0.05}, \code{miss-budget 6}, \code{giveup 0},
\code{gen-bias 0.9}, \code{p\_twist 0.0016} (neg-only,
\code{twist-min-len 256}), windows $[2,5]$ with largest-first prefix
descent, mixed sampler, ctrl-cap 2, degree $\leq 9$, span $\leq 30$, term
caps 1024/2048, \code{stop-frac 0.05}, seed 7. Question: \textbf{how many
gates does phase A end with, as a function of the generation target $G$?}

\medskip
\noindent\textbf{Ratchet split rule} (children $=$ parent $+\,1$):

\begin{center}
\begin{tabular}{rrrrr}
\toprule
target $G$ & stop at & final gates & inflation & paid \\
\midrule
2 & 1.3M moves & 177{,}390 & $1.79\times$ & 84{,}464 \\
3 & 1.4M & 180{,}458 & $1.82\times$ & 87{,}747 \\
5 & 1.5M & 182{,}742 & $1.85\times$ & 90{,}307 \\
10 & 1.7M & 186{,}632 & $1.88\times$ & 94{,}984 \\
\bottomrule
\end{tabular}
\end{center}

\noindent\textbf{Inherit split rule} (\code{-{}-gen-split-inherit}; only DB
replacements raise generations):

\begin{center}
\begin{tabular}{rrrrr}
\toprule
target $G$ & stop at & final gates & inflation & paid \\
\midrule
2 & 1.4M moves & 181{,}377 & $1.83\times$ & 88{,}807 \\
3 & 1.4M & 180{,}139 & $1.82\times$ & 87{,}553 \\
5 & 1.6M & 185{,}783 & $1.88\times$ & 94{,}263 \\
10 & 1.7M & 186{,}391 & $1.88\times$ & 95{,}423 \\
20 & 2.0M & 193{,}405 & $1.95\times$ & 104{,}612 \\
30 & 2.3M & 199{,}288 & $2.01\times$ & 112{,}878 \\
40 & 2.7M & 210{,}934 & $2.13\times$ & 129{,}236 \\
50 & 3.1M & 221{,}234 & $2.23\times$ & 142{,}279 \\
\bottomrule
\end{tabular}
\end{center}

\noindent\textbf{Lower-median variant} (\code{-{}-gen-median-low}, inherit,
$G=5$): 188{,}127 gates at 1.7M moves --- versus 185{,}783 at 1.6M for the
upper median.

\section{Interpretation}
\begin{enumerate}
\item \textbf{The first generation is the whole entry price.} Reaching
$G=2$ costs $\sim 1.8\times$ --- the coupon-collector sweep that touches
every position once. The \code{paid=} ledger (85--95k gates in every run)
accounts for essentially all growth: what phase A buys with size is the
base coverage of the material that has no free spelling, once.
\item \textbf{The curve is flat through $G\approx 10$, then gently
linear.} From $G=10$ to $G=50$ the cost rises by a steady $\sim$870 gates
($\sim$0.9\% of the input) and $\sim$35k moves per additional level:
$1.88\times$ at $G=10$ $\to$ $2.23\times$ at $G=50$. The mid-run shape
shows why: the bulk of the circuit rides the median rule's compounding (one
re-splice lifts a whole neighborhood a level), while the linear tail is the
targeted machinery grinding the bottom 5\% level by level at the splice
cadence.
\item \textbf{The stamp-rule choices barely matter --- because targeting
neutralizes them.} Ratchet vs.\ inherit: differences at the noise level
(the split ratchet is a minor accelerant; the median rule dominates). Upper
vs.\ lower median (the min on 2-gate windows, the most common splice):
$\sim$1\% in size, $\sim$100k moves. Whatever gates a stingier rule leaves
behind simply \emph{become} the laggard seed pool and get ground down at
the same cadence. The cost curve is set by the targeting + dose-stop
machinery; the stamp semantics only choose where the flat regime ends.
(The one semantics that would change the picture qualitatively is
min-over-window: every straggler drags its whole window down, and the cost
of $G$ becomes $\sim$linear from the start.)
\item \textbf{Practical consequence: choose $G$ by security need, not
budget.} The marginal cost of generation depth is so low beyond the first
sweep that $G=20$--$30$ is affordable insurance. The right value should
come from the still-open dose-response question --- how degree-2
reconstruction of the original computation decays with re-encoding depth
--- not from size pressure.
\item \textbf{Twist coverage rides along.} The benchmark runs accumulated
$\mathrm{cov}=178$--$374\times$ incidentally; the $\sim 600\times$
saturation target needs either the longer high-$G$ runs or a modestly
raised \code{p\_twist}, and can be armed as a joint stopping condition
(\code{-{}-twist-cov-stop}).
\end{enumerate}

\section{Open threads}
\begin{itemize}
\item \textbf{Generation $\leftrightarrow$ leakage dose-response}: measure
degree-2 reconstruction of the original circuit against generation depth at
a tractable size; use it to pick the production $G$.
\item \textbf{Dose-then-shrink}: after the dose stop, a compress-heavy
epoch (\code{w\_db 1}, \code{p\_db 0}) to bring the size back toward the
phase-B switch size. Generations survive compression by construction
(merges take the min, compressing splices still stamp median$+1$).
\item \textbf{Pay-ordering refinement}: hold the paid channel until the
cheap tier is nearly empty, and/or let cheap rounds take first crack at
paid products --- both would shave the $\sim G$-fold re-payment of hard
cores.
\item \textbf{Static positional-uniformity instrument} (sliding-window gate
statistics along the circuit) --- the other distinguisher front, still to
be built.
\end{itemize}
\end{document}
